\documentclass[tikz,border=1.35mm]{standalone}

\definecolor{fvmadaptedge}{HTML}{8A7E7E}
\usetikzlibrary{arrows.meta}

\definecolor{splitred}{HTML}{AF2525}

% Wireframe exploded view of prism split code 1. The top and bottom child
% prisms duplicate the new triangular split face when separated.
\providecommand{\fvmadaptYawAngle}{8}
\providecommand{\fvmadaptExplodeAmount}{1}
\providecommand{\fvmadaptExplodeScale}{1.25}
\pgfmathsetmacro{\yawcos}{cos(\fvmadaptYawAngle)}
\pgfmathsetmacro{\yawsin}{sin(\fvmadaptYawAngle)}
\pgfmathsetmacro{\xbasisx}{20.3*\yawcos + 15.3*\yawsin}
\pgfmathsetmacro{\xbasisy}{-8.5*\yawcos + 11.3*\yawsin}
\pgfmathsetmacro{\ybasisx}{-20.3*\yawsin + 15.3*\yawcos}
\pgfmathsetmacro{\ybasisy}{8.5*\yawsin + 11.3*\yawcos}
\pgfmathsetmacro{\fvmadaptEtaProjectionScale}{0.75}
\pgfmathsetmacro{\ybasisxscaled}{\fvmadaptEtaProjectionScale*\ybasisx}
\pgfmathsetmacro{\ybasisyscaled}{\fvmadaptEtaProjectionScale*\ybasisy}
\pgfmathsetmacro{\explodeZeta}{0.74*\fvmadaptExplodeScale*\fvmadaptExplodeAmount}

\begin{document}
\begin{tikzpicture}[
  x={(\xbasisx mm,\xbasisy mm)},
  y={(\ybasisxscaled mm,\ybasisyscaled mm)},
  z={(0mm,20mm)},
  line join=round,
  line cap=round,
  outer/.style={draw=fvmadaptedge,line width=0.54mm},
  split/.style={draw=splitred,line width=0.44mm},
  axis/.style={draw=fvmadaptedge!60,line width=0.32mm,-{Latex[length=2.5mm,width=1.8mm]}}
]
  % Fixed canvas for static and animated renders.
  \path[use as bounding box]
    (canvas cs:x=-50mm,y=-57mm) rectangle (canvas cs:x=58mm,y=68mm);

  \newcommand{\drawloweraxialchild}{%
    \draw[outer] (A) -- (B) -- (C) -- cycle;
    \draw[outer] (A) -- (D);
    \draw[outer] (B) -- (E);
    \draw[outer] (C) -- (F);
    \draw[split] (D) -- (E) -- (F) -- cycle;
    \foreach \p in {A,B,C} {
      \fill[fvmadaptedge] (\p) circle[radius=0.58mm];
    }
    \foreach \p in {D,E,F} {
      \fill[splitred] (\p) circle[radius=0.52mm];
    }
  }

  \newcommand{\drawupperaxialchild}{%
    \draw[split] (A) -- (B) -- (C) -- cycle;
    \draw[outer] (A) -- (D);
    \draw[outer] (B) -- (E);
    \draw[outer] (C) -- (F);
    \draw[outer] (D) -- (E) -- (F) -- cycle;
    \foreach \p in {A,B,C} {
      \fill[splitred] (\p) circle[radius=0.52mm];
    }
    \foreach \p in {D,E,F} {
      \fill[fvmadaptedge] (\p) circle[radius=0.58mm];
    }
  }

  % Compact orientation axes near the fixed lower child.
  \coordinate (Oaxis) at (-1.28,-0.55,-0.68);
  \draw[axis] (Oaxis) -- (-0.32,-0.55,-0.68) node[below right,font=\normalsize,text=fvmadaptedge] {$\xi$};
  \draw[axis] (Oaxis) -- (-1.28,0.39,-0.68) node[above left,font=\normalsize,text=fvmadaptedge] {$\eta$};
  \draw[axis] (Oaxis) -- (-1.28,-0.55,0.28) node[above,font=\normalsize,text=fvmadaptedge] {$\zeta$};

  % Lower child stays fixed.
  \begin{scope}[shift={(0,0,0)}]
    \coordinate (A) at (0,0,0);
    \coordinate (B) at (1.45,0,0);
    \coordinate (C) at (0,1,0);
    \coordinate (D) at (0,0,0.5);
    \coordinate (E) at (1.45,0,0.5);
    \coordinate (F) at (0,1,0.5);
    \drawloweraxialchild
  \end{scope}

  % Upper child lifts away along zeta.
  \begin{scope}[shift={(0,0,\explodeZeta)}]
    \coordinate (A) at (0,0,0.5);
    \coordinate (B) at (1.45,0,0.5);
    \coordinate (C) at (0,1,0.5);
    \coordinate (D) at (0,0,1);
    \coordinate (E) at (1.45,0,1);
    \coordinate (F) at (0,1,1);
    \drawupperaxialchild
  \end{scope}
\end{tikzpicture}
\end{document}
